La metrología es una disciplina fundamental en la ingeniería, considerada el arte de medir. Su objetivo principal es garantizar que las mediciones sean precisas y confiables, permitiendo que los resultados obtenidos sean comparables y verificables. En este contexto, el diseño de instrumentos de medición precisos y fiables es de crucial importancia, y estos deben ser contrastados con otros instrumentos de mayor precisión para asegurar la validez de las mediciones.

Es igualmente importante comprender el impacto que tiene la medición sobre el propio sistema. Al incorporar un instrumento de medición a un circuito o sistema, se introduce una perturbación inevitable, ya que el propio dispositivo de medición altera las condiciones del sistema. Esto implica que nunca será posible conocer los valores exactos de las magnitudes que se miden. Sin embargo, mediante un análisis adecuado, es posible determinar una cota de error que nos permita aproximarnos a los valores reales, garantizando la confiabilidad de los resultados.

En la primera parte de esta experiencia de laboratorio, se exploró el funcionamiento del galvanómetro y su comportamiento como filtro pasa bajos. También se realizó el diseño y prueba de un voltímetro analógico, utilizando este mismo galvanómetro. Como continuación natural de este trabajo, en esta segunda parte se procederá a la construcción de un amperímetro analógico, también basado en el galvanómetro de bobina móvil.

Para la realización de este amperímetro, se llevarán a cabo cálculos teóricos, simulaciones, y finalmente el diseño físico del instrumento. Este se calibrará adecuadamente y se pondrá a prueba en distintos rangos de medición, contrastando sus resultados con los obtenidos por instrumentos de mayor precisión para verificar su exactitud.

\subsection{Objetivos}
\subsubsection*{Objetivos Generales}
\begin{itemize}
    \item Diseñar un amperímetro analógico de tres rangos a partir de un galvanómetro de bobina móvil.
    \item Comprender el principio de funcionamiento de un amperímetro analógico y las protecciones que debe incluir. 
\end{itemize}

\subsubsection*{Objetivos Específicos}
\begin{itemize}
    \item Calibrar un amperímetro para el máximo rango de medición
    \item Medir el error inducido en mediciones de un rango inferior 
    \item Comprobar la linealidad de un instrumento analógico compuesto por elementos pasivos.
\end{itemize}